% Figura do slide 04. Compilar a partir desta pasta.
\documentclass[aspectratio=169,11pt]{beamer}
% Estilo das figuras: mesmas fontes, tamanhos e cores da apresentação.
\usepackage[utf8]{inputenc}
\usepackage[T1]{fontenc}
\usepackage[brazilian]{babel}
\usepackage{amsmath,tikz}
\usetikzlibrary{arrows.meta,positioning,calc}
\usepackage[active,tightpage]{preview}
\PreviewEnvironment{tikzpicture}
\setlength{\PreviewBorder}{0pt}
\definecolor{ink}{HTML}{202124}
\definecolor{muted}{HTML}{666666}
\definecolor{blue}{RGB}{30,80,160}
\definecolor{green}{RGB}{30,130,60}
\definecolor{red}{RGB}{190,60,60}
\definecolor{light}{HTML}{E8E8E8}
\definecolor{amber}{HTML}{B68432}
\setbeamercolor{normal text}{fg=ink,bg=white}
\setbeamersize{text margin left=8mm,text margin right=8mm}
\setbeamertemplate{navigation symbols}{}
\tikzset{arr/.style={-{Stealth[length=2mm]},thick,muted},txt/.style={align=center,font=\normalsize},smalltxt/.style={align=center,font=\small,text=muted}}

\begin{document}
\begin{frame}
\begin{tikzpicture}
\node[anchor=west,font=\large] at (0,0) {Território e continuidade};
\node[anchor=west,smalltxt] at (0,-.55) {Faria (2020) · CONASS (2021)};
\node[anchor=west,font=\large] at (8.3,0) {Polígonos e prazos};
\draw[arr] (6.7,0)--(7.8,0);
\node[anchor=west,font=\large] at (0,-1.55) {Informação e gestão};
\node[anchor=west,smalltxt] at (0,-2.1) {Rodrigues et al. (2014) · Bender et al. (2024)};
\node[anchor=west,font=\large] at (8.3,-1.55) {Histórico e exportação};
\draw[arr] (6.7,-1.55)--(7.8,-1.55);
\node[anchor=west,font=\large] at (0,-3.1) {Roteamento em saúde};
\node[anchor=west,smalltxt] at (0,-3.65) {Randriamihaja et al. (2024)};
\node[anchor=west,font=\large] at (8.3,-3.1) {Rotas com jornada};
\draw[arr] (6.7,-3.1)--(7.8,-3.1);
\end{tikzpicture}
\end{frame}
\end{document}
